\chapter{Introduction}
Flying-wing UAS design forces you to confront coupling early. Planform and twist choices affect trim and static margin. Stability choices affect control surface deflections and spanwise loading, which changes structural margins and the aerodynamic drag picture. Propulsion and energy storage choices change mass properties and mission feasibility, and those changes feed back into the aero and structure.

This literature review focuses on methods and tools that directly support an integrated preliminary design workflow for flying-wing aircraft. The emphasis is on what the sources actually do, what they assume, and where they tend to fail. The goal is to extract guidance for a tool that connects geometry definition, aerodynamic analysis, mission simulation, structural checks, and export to downstream environments.

\chapter{Summary of Reviewed Sources}
\begingroup
\small
\setlength{\tabcolsep}{6pt}
\renewcommand{\arraystretch}{1.25}
\begin{longtable}{p{0.16\linewidth} p{0.23\linewidth} p{0.26\linewidth} p{0.23\linewidth}}
\toprule
\textbf{Source} & \textbf{What it does} & \textbf{Key details used in this review} & \textbf{Access notes} \\
\midrule
\endfirsthead
\toprule
\textbf{Source} & \textbf{What it does} & \textbf{Key details used in this review} & \textbf{Access notes} \\
\midrule
\endhead
\midrule
\multicolumn{4}{r}{\textit{Continued on next page}} \\
\midrule
\endfoot
\bottomrule
\endlastfoot

Nguyen et al. (2013) \cite{Nguyen2013MFM} &
Multi-fidelity MDO for a UCAV concept using a CFD-backed metamodel inside an MDF loop. &
D-optimal DOE for CFD points; CFD cruise condition reported as Mach 0.85 at 38{,}000 ft and 2 deg; DOE size reported as 20 configs; metamodel used to avoid CFD-in-the-loop. &
Often paywalled (journal); look for author copy or institutional access. \\

Vegh et al. (2019) \cite{Vegh2019NDARCSUAVE} &
Cross-tool comparison (NDARC vs SUAVE) on Kitty Hawk Cora with sweeps and SLSQP optimization. &
Baseline GTOW reported (2233 lb vs 2086 lb); wingspan sweep trend disagreement attributed to drag; optimized designs and cross-evaluation reported. &
Public PDF available via SUAVE/Stanford. \\

Deperrois (2009) \cite{Deperrois2009XFLR5} &
XFLR5 guidelines describing scope, assumptions, and numerical setup expectations. &
Explicit warnings about applicability; emphasizes mesh/panel refinement and convergence checks. &
Public PDF (XFLR5 documentation). \\

Yu et al. (2022) \cite{Yu2022ICAS} &
Potential-flow solver comparison on T-FLEX UAV against STAR-CCM+ Euler/RANS baselines. &
RANS uses Spalart-Allmaras; CFD mesh study 0.9--22M cells (12.3M chosen); potential-flow convergence criterion (1\% CL change); notes on $C_m$ and induced drag extraction differences. &
Public PDF (ICAS archive). \\

Rosas-Cordova et al. (2024) \cite{RosasCordova2024VSPAERO} &
VSPAERO validation and tessellation sensitivity on basic wings. &
Spanwise tessellation drives convergence; recommends spanwise tessellation $>4$ for unswept; reports lift-curve-slope errors (20.55\% at AR=2, 6.19\% at AR=4); provides a recommended tessellation parameter set and stated validity bounds. &
Public PDF available via the publisher site mirror. \\

Zhang et al. (2025) \cite{Zhang2025ANNGA} &
Flying-wing glider optimization using ANN surrogate and GA with stability constraints. &
Static margin constraint reported as 5\%; reports performance penalty (11.5\% range reduction) and example stabilized performance numbers. &
MDPI article is typically open-access. \\

Sharpe (2021) \cite{Sharpe2021AeroSandbox} &
AeroSandbox thesis: differentiable analysis and optimization-first workflows for aircraft design. &
Discussion of ``glass-box'' requirements under IPOPT; framing of SAND-style formulations; emphasis on robust analysis functions for optimization. &
MIT thesis repository (sometimes blocked by network policies); may need manual download. \\

Tremblay et al. (2007) \cite{Tremblay2007Battery} &
Generic Thevenin-style battery model for dynamic simulation with parameter identification from discharge curves. &
Terminal-voltage model structure (OCV + polarization + exponential zone + ohmic drop); stated omissions (no hysteresis, no temperature dependence). &
IEEE paper (often paywalled). \\

Berndt (2009) \cite{Berndt2009JSBSim} &
JSBSim flight dynamics model description and intended modeling scope. &
Used to justify export-to-simulator as a verification path and to frame what is reasonable to expect from the external model. &
AIAA paper (often paywalled). \\

Bruhn; Niu; NASA SP-8007; Timoshenko and Gere \cite{Bruhn,Niu,NASASP8007,TimoshenkoGere} &
Standard references for aircraft structural sizing and buckling checks. &
Used to motivate the checklist of early structural outputs: stress, deflection, torsion, panel buckling, column/crippling margins. &
Books (library access) and NASA report (public). \\

\end{longtable}
\endgroup

\chapter{Multi-Fidelity Preliminary Design and System-Level Coupling}
\section{Multi-fidelity MDO for conceptual design}
Nguyen et al. describe a multidisciplinary conceptual design approach for an unmanned combat air vehicle that explicitly mixes low-fidelity discipline models with a higher-fidelity aerodynamic correction strategy \cite{Nguyen2013MFM}. Their baseline conceptual loop couples aerodynamic estimation, weight estimation, and mission range inside a multidisciplinary feasible (MDF) iteration. Instead of embedding CFD inside the optimizer, they generate a limited set of higher-fidelity aerodynamic evaluations and fit a metamodel that the MDO process can query quickly.

Two technical choices in Nguyen et al. map well to a tool workflow. First, they do not sample high-fidelity points randomly. They use D-optimal design of experiments to choose configurations that improve metamodel quality per CFD run \cite{Nguyen2013MFM}. Second, they evaluate those CFD points at an explicitly stated cruise condition (reported as Mach 0.85, 38,000 ft, and 2 deg angle of attack) using Fluent, then use the metamodel in the optimizer \cite{Nguyen2013MFM}. The high-fidelity effort is used to correct a low-order model in a narrow but important operating regime.

They report improved aerodynamic prediction accuracy without a commensurate increase in design time because the optimizer never directly pays CFD cost \cite{Nguyen2013MFM}. The limitation is also structural: as the design moves away from the DOE region, the metamodel becomes an extrapolator. A robust workflow therefore needs both diagnostics (to detect leaving the training envelope) and a mechanism to add new high-fidelity points when needed.

In their implementation, the DOE is small by design (reported as 20 configurations for the high-fidelity database) and the optimization is posed as a multi-objective trade between mission range and aircraft weight \cite{Nguyen2013MFM}. The method is useful because it combines fast iteration with a higher-fidelity anchor while making the validity region explicit. The tool is only as trustworthy as the DOE coverage.

\section{Cross-tool comparison as verification}
Vegh et al. compare NDARC and SUAVE by modeling the Kitty Hawk Cora aircraft under a simplified mission definition \cite{Vegh2019NDARCSUAVE}. The aircraft class differs from a flying wing, but their main contribution here is methodological: it shows how two established conceptual design tools can disagree in ways that change design decisions.

They report baseline gross takeoff weight (GTOW) predictions in the same range (2233 lb for NDARC and 2086 lb for SUAVE) while noting that the internal modeling choices differ in meaningful ways \cite{Vegh2019NDARCSUAVE}. For example, SUAVE performs a mission simulation that converges an energy network at each time step and represents the battery using a Thevenin-equivalent circuit, whereas NDARC uses a more sizing-oriented approach \cite{Vegh2019NDARCSUAVE}. When they sweep rotor radius, both tools show similar weight sensitivity. When they sweep wingspan, the tools show opposite trends, which they attribute primarily to differences in wing drag predictions \cite{Vegh2019NDARCSUAVE}.

The paper is also careful about defining the scenario it is comparing against. They collect public data for the Cora concept (including an 11 m wingspan) and use a simplified Uber Elevate mission definition with a representative payload (reported as 400 lb), a range target (100 km), and a cruise speed target (180 km/h), along with an operating altitude band (500 to 3000 ft) \cite{Vegh2019NDARCSUAVE}. Those details matter because many disagreements between conceptual tools appear to be physics disagreements but are actually mismatches in mission definition and accounting conventions.

The paper also reports an optimization study using SLSQP. The resulting optima differ not only in the chosen design variables, but also in GTOW. They then cross-evaluate each optimized design in the other tool and show that the discrepancy persists \cite{Vegh2019NDARCSUAVE}. The takeaway for an integrated preliminary design tool is that cross-checking is not a luxury. A tool should make it easy to reproduce intermediate quantities (lift, drag breakdowns, power flows, margins) so disagreements can be traced to assumptions, conventions, or numerical setup rather than being waved away.

One detail that makes the comparison actionable is that they report the optimized designs, not just the fact that they differ. In their study, the NDARC optimum is reported at 2223 lb with a rotor radius of 1.83 ft and wingspan of 39.5 ft, while the SUAVE optimum is reported at 1948 lb with a rotor radius of 1.70 ft and wingspan of 30.0 ft \cite{Vegh2019NDARCSUAVE}. Even if the exact numbers are not transferable to a flying wing, the pattern is: two tools can produce different optima under the same requirements because they disagree on intermediate modeling components like drag and energy usage.

\chapter{Aerodynamic Analysis for UAS and Flying Wings}
\section{What XFLR5 is designed to do (and what it is not)}
Deperrois' XFLR5 guidelines are direct about scope and limitations \cite{Deperrois2009XFLR5}. The document positions XFLR5 as a tool for model sailplanes and model aircraft and explicitly discourages using it for full-size aircraft design \cite{Deperrois2009XFLR5}. The software couples 2D airfoil analysis (via XFOIL-style methods) with 3D wing analysis using LLT, VLM, and a panel method. Viscous effects enter primarily through section polars rather than a full 3D viscous solution.

The guidelines emphasize that numerical choices matter: panel discretization and refinement influence results, and users are expected to run convergence checks \cite{Deperrois2009XFLR5}. For a flying-wing workflow, XFLR5 is useful as a fast trend tool when its assumptions hold. When designs approach stall, strong separation, or complex control-surface behavior, its outputs are best treated as hypotheses to test rather than answers to trust.

\section{Potential-flow solver comparison on a UAV configuration}
Yu et al. compare a set of potential-flow solvers on the T-FLEX UAV and evaluate the results against higher-fidelity CFD performed in STAR-CCM+ \cite{Yu2022ICAS}. They include both Euler and RANS simulations, and for the RANS baseline they report using the Spalart-Allmaras turbulence model. They document a mesh study from 0.9 to 22 million cells and select a 12.3 million cell mesh as a convergence compromise \cite{Yu2022ICAS}.

The set of potential-flow tools they evaluate includes AVL, Tornado, XFLR5, VSPAERO, PAWAT, and FlightStream \cite{Yu2022ICAS}. What the paper adds is a clear comparison of how those tools differ on the same configuration, rather than treating them as interchangeable.

A strength of the study is that they treat numerical convergence as part of the experiment. On the potential-flow side, they increase panel resolution until the change in lift coefficient between refinements is below 1\% \cite{Yu2022ICAS}. This matters because solver-to-solver differences can be smaller than the error induced by an under-resolved discretization.

Their findings are nuanced. In the linear regime, lift curve trends are generally consistent across tools, but drag and moment behavior is more sensitive to model structure and post-processing. They discuss that some tools incorporate boundary-layer corrections (FlightStream is described as using an integral boundary layer method), while others remain closer to inviscid potential flow \cite{Yu2022ICAS}. For pitching moment results, they report that VSPAERO and FlightStream agree more closely with the CFD pitching moment slope for their configuration \cite{Yu2022ICAS}. They also note that induced drag is not a single, universally extracted quantity in these codes; how induced drag is computed and reported can drive apparent disagreement even when lift predictions match \cite{Yu2022ICAS}.

For flying-wing design, this distinction is not academic. A design can look excellent in lift and induced drag and still be unusable because the pitching moment curve forces large trim deflections or produces an unstable neutral point. Yu et al.'s comparison implicitly ranks outputs: for conceptual trade studies, a solver that provides stable and repeatable $C_m(\\alpha)$ behavior is often more valuable than one that only provides $C_L$ trends \cite{Yu2022ICAS}.

For a flying-wing preliminary design tool, these details matter because the design is often constrained by moment and stability behavior at least as much as by L/D. A solver backend that provides lift but gives unreliable moment trends can mislead an early design loop.

\section{VSPAERO validation and mesh recommendations}
Rosas-Cordova et al. validate VSPAERO on a set of basic wing configurations and focus on how tessellation choices affect results \cite{RosasCordova2024VSPAERO}. They provide mesh parameter recommendations and quantify where the method struggles.

They report that chordwise tessellation had little effect on their results in the studied cases compared with other uncertainties, while spanwise tessellation strongly affected convergence. For unswept wings they recommend spanwise tessellation values above 4 for consistent convergence, and they note that swept wings require finer spanwise resolution \cite{RosasCordova2024VSPAERO}. They also summarize bounds where they consider results acceptable (for example, moderate angles of attack and higher aspect ratio cases) and explicitly call out regimes where errors grow \cite{RosasCordova2024VSPAERO}.

They provide error magnitudes that are relevant to tool users. In a low Reynolds number case with aspect ratio 2, they report a 20.55\% error in lift curve slope prediction, while a higher aspect ratio case (AR=4) shows a smaller error (6.19\%) \cite{RosasCordova2024VSPAERO}. Flying wings often sit in the moderate-to-low aspect ratio regime, so this result should shape expectations. VSPAERO can be a fast baseline for planform trade studies, but results should be cross-checked for low-AR wings and for swept cases where tip stall and 3D separation effects become important.

They also compare VSPAERO outputs to experimental references (including classic NACA data for basic wing cases) and summarize where the method tends to under- or over-predict trends \cite{RosasCordova2024VSPAERO}. They make the point that validation is rarely one number. A tool can match lift curve slope while missing pitching moment slope, or match moment while missing drag. Their paper is most useful where it separates numerical setup effects (tessellation choices) from modeling-domain effects (aspect ratio, Reynolds number, sweep) \cite{RosasCordova2024VSPAERO}.

They also provide a concrete starting point for meshing parameters that can be embedded in a tool as defaults. One recommended set includes chordwise tessellation of 33, spanwise tessellation of 24, chordwise cluster of 0.3, spanwise cluster of 0.6, wake length of 5 chord lengths, and wake tessellation of 24 \cite{RosasCordova2024VSPAERO}. They describe an intended validity regime that includes moderate angles of attack (reported as -10 deg to 10 deg), low Mach number (M < 0.3), Reynolds number above 100,000, and aspect ratio above 4 \cite{RosasCordova2024VSPAERO}. Those bounds are useful because they separate two questions a design tool user might otherwise conflate: whether the solver ran and whether the solver is expected to be accurate.

\chapter{Flying-Wing Configuration, Stability, and Design Optimization}
\section{Optimization with explicit stability constraints}
Zhang et al. present an optimization workflow for flying-wing gliders that treats stability as a constraint rather than a post-check \cite{Zhang2025ANNGA}. Their approach uses an ANN-based surrogate model coupled with a genetic algorithm. They include a static margin requirement (reported as 5\%) and evaluate the performance cost of enforcing it \cite{Zhang2025ANNGA}.

Methodologically, the paper is an example of ``surrogate-first'' preliminary design. The ANN stands in for repeated aerodynamic evaluations inside the genetic algorithm loop, which allows broader exploration of the design space than would be practical if each candidate required a slow solver call \cite{Zhang2025ANNGA}. The tradeoff is that the optimizer inherits the surrogate's blind spots; stability and trim constraints therefore do double duty as both design requirements and as guardrails that reduce the chance of selecting an obviously unrealistic configuration.

Their results make two points that align with practical flying-wing design experience. First, enforcing stability has a measurable performance penalty. They report an 11.5\% reduction in maximum range relative to an unstable optimum, while still achieving a reported glide ratio of 24.1 and a static margin of 5.1\% \cite{Zhang2025ANNGA}. Second, stability enforcement can drive geometry changes that are not obvious from aerodynamic efficiency alone, because the available degrees of freedom are limited when there is no tail volume.

They also illustrate a limitation common to surrogate-driven optimization: the surrogate is only as good as the training data and the aerodynamic model used to generate it. If the training method does not capture separation onset, control-surface effectiveness near stall, or post-stall moment behavior, the optimizer can select designs that look attractive in the surrogate but fail in higher-fidelity analysis.

\chapter{Differentiable Analysis and Optimization-Oriented Tooling}
\section{AeroSandbox and glass-box optimization requirements}
Sharpe's AeroSandbox thesis is explicit about building aircraft analysis around optimization rather than adding optimization later \cite{Sharpe2021AeroSandbox}. The implementation is built around algorithmic differentiation and nonlinear programming, using CasADi for differentiation and IPOPT as a primary solver \cite{Sharpe2021AeroSandbox}.

A central practical issue discussed in the thesis is that IPOPT expects analysis functions to behave like ``glass-box'' mappings. Analysis functions must return outputs and gradients in a form that remains defined even when intermediate values wander into infeasible regions. This motivates robust extrapolation behavior, careful variable bounds, and constraint formulations that do not produce undefined expressions \cite{Sharpe2021AeroSandbox}.

The thesis also motivates simultaneous analysis and design (SAND) formulations: instead of nesting analysis inside optimization, analysis state variables can be included as decision variables with constraints enforcing the governing equations \cite{Sharpe2021AeroSandbox}. In preliminary design contexts, SAND can reduce outer-loop iteration counts and make gradients more informative, but it increases the size of the nonlinear program and places more burden on numerical conditioning.

For a flying-wing workflow, the mapping is direct. Twist optimization, lift distribution shaping, and mission parameter tuning can be posed as constrained optimizations, but only if the underlying analysis functions stay well-behaved under optimizer pressure. Sharpe is explicit about where the effort goes if you want the optimization results to be trustworthy: not the optimizer call, but the analysis robustness \cite{Sharpe2021AeroSandbox}.

\chapter{Mission Modeling, Interoperability, and Export}
Mission modeling often fails first at the interfaces. Vegh et al. show this indirectly: even when two tools target the same vehicle and similar mission requirements, differences in aerodynamic drag modeling and mission simulation details can flip trade study trends \cite{Vegh2019NDARCSUAVE}. For an integrated design tool, this suggests that mission simulation outputs should preserve intermediate quantities (forces, power, state-of-charge trajectories, margins) so the user can see why the answer changes.

If the workflow exports models to an external simulator, the internal modeling conventions must map cleanly to the simulator's configuration language. Berndt describes JSBSim as an open-source flight dynamics model and clarifies its intended modeling scope \cite{Berndt2009JSBSim}. That makes it a reasonable target for early verification of trim and stability behavior. The export path then becomes part of the modeling contract: coordinate frames, sign conventions, reference areas, and control surface definitions must be explicit, consistent, and testable.

\chapter{Propulsion and Energy Storage Models}
\section{A battery model that is simple but parameterizable}
Tremblay et al. propose a generic battery model intended for dynamic simulation, built around a small set of parameters that can be identified from discharge curves \cite{Tremblay2007Battery}. The model represents the terminal voltage as a combination of an open-circuit voltage term, a polarization term that depends on extracted capacity, an exponential term to capture initial voltage drop behavior, and a series resistance drop \cite{Tremblay2007Battery}. The state of charge is tracked through current integration, and the formulation is designed to be computationally light enough for system simulation.

In their notation, the terminal voltage is expressed as a sum of terms of the form
\[
V_{batt} = E_0 - K\frac{Q}{Q - it}\,it - K\frac{Q}{Q - it}\,i^* + A\exp(-B\,it) - R\,i,
\]
where $E_0$ is a constant voltage, $Q$ is capacity, $it$ is extracted capacity, $i$ is current, $i^*$ is a filtered current term, and $A$ and $B$ shape the exponential zone \cite{Tremblay2007Battery}. The exact algebra is less important than the behavior it captures: an initial voltage drop and a slower polarization-driven sag using a parameter set that can be fitted from a discharge curve.

They describe a practical identification procedure using manufacturer discharge data and validation by comparing simulated voltage profiles to those curves \cite{Tremblay2007Battery}. The limitations are clearly stated: the model does not include hysteresis and does not explicitly model temperature dependence, so it is best interpreted as a mission-level model rather than a cell-physics model \cite{Tremblay2007Battery}.

For electric UAS preliminary design, the value of this model class is that it exposes voltage sag and usable-energy sensitivity under power draw without requiring stiff electrochemical dynamics. That allows design trade studies to capture the coupling between propulsion sizing, mission profile, and energy margin.

\chapter{Structural Sizing in Preliminary Design}
The structural sources referenced in this review are used as engineering standards rather than as novel research contributions, but they define the baseline for what structural checks look like in early aircraft design.

Timoshenko and Gere provide core stability and elastic buckling foundations \cite{TimoshenkoGere}. Bruhn and Niu provide aircraft-oriented sizing guidance and worked examples that connect load cases to practical thickness and stiffener decisions \cite{Bruhn,Niu}. NASA SP-8007 is a standard reference for thin-walled buckling behavior and is commonly used to benchmark buckling checks \cite{NASASP8007}. In a preliminary design tool, these references support a fast check set: stress, deflection, torsion, and primary panel and column buckling limits, coupled to aerodynamic spanwise loading.

\chapter{Synthesis: Requirements for a Useful Flying-Wing Design Tool}
Across the sources, the theme is not that one method wins. It is that a useful preliminary design tool must make iteration cheap while staying honest about assumptions.

Nguyen et al. show a specific pattern for inserting higher-fidelity corrections into an MDO loop without destroying runtime: use DOE, fit a metamodel, and keep the optimizer in the corrected low-order space \cite{Nguyen2013MFM}. Yu et al. and Rosas-Cordova et al. show that solver comparisons only mean something when numerical convergence is treated as part of the setup, and that mesh parameters can dominate error if they are implicit \cite{Yu2022ICAS,RosasCordova2024VSPAERO}. Zhang et al. quantify that stability constraints are not free and can impose a real performance penalty \cite{Zhang2025ANNGA}. Sharpe provides a blueprint for optimization-first analysis, but also makes the engineering cost visible: analysis functions must be robust enough for an optimizer to push on them \cite{Sharpe2021AeroSandbox}. Tremblay et al. provide a battery model that is simple enough for mission simulation yet detailed enough to expose voltage sag and usable energy sensitivity \cite{Tremblay2007Battery}.

If the program goal is a design environment that produces actionable insights, the remaining gaps are mostly about trust and traceability: explicit conventions (signs, axes, reference areas), repeatable numerics (mesh and panel convergence policies), and automated cross-checks against at least one independent method. Without that, the tool produces fast answers that are hard to defend when the configuration changes or when the design is handed off to detailed analysis.
